% =====================================================================
% Capítulo 1. Introducción
% Estado: capítulo completo, reorganizado el 13/10/2026 según la plantilla de
% la UAX en dos apartados (justificación y presentación). El texto de las
% antiguas secciones 1.1 a 1.4 se conserva sin cambios, salvo dos referencias
% cruzadas. Párrafo de motivación personal al comienzo de la sección de
% justificación, redactado por el autor (13/10/2026).
% =====================================================================

\chapter{Introducción}\label{cap:introduccion}

Este trabajo compara distintos modelos para prever, con un día de antelación,
la demanda horaria de electricidad del sistema peninsular español. Los modelos
van desde un método ingenuo y un modelo estadístico clásico de series
temporales hasta un modelo de \textit{gradient boosting} y una red neuronal, y
se comparan bajo unas mismas reglas, de modo que las diferencias observadas
puedan atribuirse a los propios métodos. La evaluación no se limita a medir
cuánto se equivoca cada modelo: comprueba también si los intervalos de
predicción que acompañan a sus previsiones son fiables, es decir, si contienen
el valor observado con la frecuencia que anuncian.

El capítulo se organiza en dos apartados. El primero justifica el trabajo:
expone el contexto que hace relevante el problema y las carencias de la
literatura que el trabajo pretende cubrir. El segundo presenta el proyecto:
delimita el problema abordado, formula la pregunta de investigación y las
aportaciones del trabajo, y describe la organización del resto del documento.

\section{Justificación del proyecto realizado}\label{sec:justificacion}

La elección del tema responde, en primer lugar, al interés por los modelos predictivos y por la contribución de las matemáticas a la resolución de problemas reales, que desarrollé en mi formación del grado en Matemáticas. Durante este grado realicé trabajos de predicción de resultados deportivos en los que se emplearon modelos de aprendizaje automático, como los bosques aleatorios (random forest) y XGBoost, de la misma familia de gradient boosting que LightGBM, junto con modelos de regresión de Poisson y exponencial negativa. El descubrimiento posterior de las series temporales, en las que el orden de las observaciones y la dependencia entre ellas forman parte del propio problema, despertó mi interés por trabajar con este tipo de datos. Además, el sector energético se me presentó como el ámbito más adecuado para hacerlo, por tratarse de un sector fundamental en el presente y en el futuro del planeta.

La justificación del trabajo se apoya en dos argumentos, que se desarrollan a
continuación. El primero es la relevancia práctica del problema: la previsión
de la demanda condiciona la operación del sistema eléctrico y la participación
en el mercado, y las transformaciones recientes del sector han aumentado su
exigencia. El segundo es su interés científico: la literatura que compara
modelos estadísticos y de aprendizaje automático presenta carencias que dejan
abierta la cuestión de cuál de ellos conviene emplear en este problema.

\subsection{Contexto y relevancia del problema}\label{sec:contexto}

El sistema eléctrico presenta una particularidad que lo distingue de la mayoría
de los mercados de bienes: la producción y el consumo deben igualarse en cada
instante. Aunque la capacidad de almacenamiento está creciendo, sigue siendo
reducida en relación con el volumen de demanda que el sistema atiende, de modo
que la energía consumida ha de generarse prácticamente en el mismo momento en
que se utiliza. Esta restricción física sitúa la anticipación de la demanda por
delante de cualquier decisión de operación: qué unidades de generación se
programan, cuánta reserva se mantiene disponible y a qué precio se intercambia
la energía dependen, en último término, de una estimación de cuánta
electricidad se va a consumir.

La previsión de la demanda se elabora a distintos horizontes en función de la
decisión a la que sirve. La planificación de la red y de la capacidad de
generación requiere previsiones a varios años vista, mientras que la operación
del sistema y la participación en los mercados se apoyan en previsiones que
abarcan desde unas pocas horas hasta algunos días \parencite{hong2016probabilistic}.
Este trabajo se sitúa en el segundo grupo, la previsión a corto plazo, cuyas
consecuencias económicas se materializan de forma inmediata.

En España, la Ley~24/2013, del Sector Eléctrico, atribuye a Red Eléctrica de
España la función de operador del sistema, cuyo cometido principal es
garantizar la continuidad y la seguridad del suministro y coordinar
adecuadamente la producción y el transporte de energía \parencite{ley2013sectorelectrico}.
La gestión del mercado organizado de electricidad corresponde, por su parte, al
operador del mercado, función que desempeña el Operador del Mercado Ibérico de
Energía (OMIE).

En el mercado diario, los agentes presentan sus ofertas de compra y venta de
energía para el día siguiente, y la recepción de ofertas se cierra a las
12:00, hora peninsular \parencite{omie2025mtu15}. Tanto la comercializadora que adquiere
energía para sus clientes como el generador que decide cuánta ofrecer lo hacen,
por tanto, a partir de una previsión de la demanda de las veinticuatro horas
siguientes elaborada antes de ese cierre. Tras la casación, el operador del
sistema resuelve las restricciones técnicas y gestiona los servicios de ajuste
necesarios para equilibrar producción y consumo en tiempo real. Las diferencias
entre la energía programada y la finalmente consumida, conocidas como desvíos,
se liquidan y generan costes para los agentes responsables. Un error de
previsión se traduce así en compras de energía mal dimensionadas, en una mayor
necesidad de reservas y en costes de desvío.

El marco de funcionamiento, además, se ha vuelto más exigente en los últimos
años. Desde el 1 de diciembre de 2024, los desvíos se liquidan en periodos de
quince minutos \parencite{cnmc2025procedimientos}, y desde la jornada de entrega
del 1 de octubre de 2025 el mercado diario casa productos cuartohorarios, con
96 precios diarios en lugar de 24 \parencite{omie2025mtu15}. Al reducirse la
unidad temporal en la que se compensan los errores, la precisión de la
previsión gana peso económico. Este trabajo se circunscribe al periodo
2017-2024, durante el cual la casación del mercado diario, que es la decisión
a la que sirve la previsión estudiada, se realizó íntegramente en resolución
horaria. La delimitación completa del problema se presenta en la
sección~\ref{sec:problema}.

A ello se suma la transformación del propio sistema. La creciente participación
de fuentes renovables no gestionables, como la eólica y la solar fotovoltaica,
obliga a equilibrar una generación variable con una demanda que también lo es.
La expansión del autoconsumo, por su parte, modifica la forma de la curva de
demanda que observa el sistema, especialmente en las horas centrales del día.
Ambos fenómenos aumentan el valor de disponer de previsiones fiables y, al mismo
tiempo, hacen menos probable que los patrones históricos se mantengan estables.

Por último, una previsión puntual, es decir, un único valor para cada hora,
resulta insuficiente para buena parte de estas decisiones. Dimensionar las
reservas o gestionar el riesgo de desvío exige conocer también la incertidumbre
asociada a la previsión, lo que ha impulsado el desarrollo de la previsión
probabilística de la demanda \parencite{hong2016probabilistic}. Un intervalo de
predicción, sin embargo, solo es útil si está calibrado: si un intervalo
nominal del 90\,\% contiene el valor observado con una frecuencia muy distinta
del 90\,\%, la información que transmite a quien toma la decisión es engañosa.

La confluencia de estos factores motiva el presente trabajo: un problema de
relevancia operativa y económica directa, la disponibilidad de datos públicos
de calidad publicados por el propio operador del sistema y una necesidad
creciente de cuantificar, y no solo de reducir, el error de previsión.

\subsection{Carencias de la literatura}\label{sec:carencias}

La previsión de la demanda eléctrica cuenta con una literatura extensa, en la
que conviven los modelos estadísticos clásicos con un número creciente de
propuestas basadas en aprendizaje automático y profundo. La evidencia sobre
el rendimiento relativo de ambos enfoques, sin embargo, es menos concluyente
de lo que ese volumen de publicaciones podría sugerir.
\textcite{makridakis2018statistical} señalaron la escasez de evidencia al
respecto y, al comparar varios métodos populares de aprendizaje automático con
ocho métodos estadísticos tradicionales sobre 1045 series mensuales de la
competición M3, encontraron que los primeros quedaban por detrás en todas las
métricas y horizontes analizados, con un coste computacional
considerablemente mayor. Aunque este resultado procede de series mensuales y
no puede trasladarse sin más a la demanda horaria, muestra que la ventaja de
los métodos más complejos no puede darse por supuesta y debe demostrarse en
cada problema.

La revisión de la literatura que se desarrolla en el marco teórico pone de
manifiesto dos limitaciones recurrentes en los estudios que comparan ambas
familias de modelos. La primera es la falta de equidad en las condiciones de
comparación. Cuando un modelo estadístico se emplea en su configuración por
defecto frente a una red neuronal cuidadosamente ajustada, o cuando los
modelos no disponen de las mismas variables explicativas ni del mismo
horizonte, las diferencias observadas no pueden atribuirse al método, porque
pueden deberse simplemente a un esfuerzo de ajuste desigual. La segunda es la
restricción de la evaluación al error puntual. La literatura de previsión de
la demanda se ha centrado mayoritariamente en la previsión puntual
\parencite{hong2016probabilistic}, a pesar de que, como se ha expuesto en la
sección~\ref{sec:contexto}, las decisiones de operación requieren conocer
también la incertidumbre asociada a la previsión.

Evaluar esa incertidumbre exige, además, comprobar que los intervalos de
predicción son fiables. El criterio de referencia consiste en maximizar la
nitidez de las previsiones probabilísticas, es decir, obtener intervalos lo
más estrechos posible, con la condición de que estén calibrados
\parencite{gneiting2007probabilistic}. La predicción conforme ofrece un
mecanismo para recalibrar los intervalos que produce cualquier modelo sin
imponer supuestos sobre la distribución de sus errores. Su aplicación a series
temporales, en las que no se cumple la intercambiabilidad de los datos que
exige la teoría conforme clásica, ha dado lugar a variantes específicas, entre
ellas la inferencia conforme adaptativa (ACI), diseñada para mantener la
cobertura a largo plazo aunque la distribución de los datos cambie con el
tiempo \parencite{gibbs2021adaptive}.

\section{Presentación del proyecto y sus contenidos}\label{sec:presentacion}

Este apartado describe el problema que aborda el trabajo y cómo contribuye a
resolverlo. En primer lugar se delimita el problema de previsión; a
continuación se formulan la pregunta de investigación y las aportaciones del
trabajo, y por último se presenta la organización del documento.

\subsection{Planteamiento y delimitación del problema}\label{sec:problema}

El objeto de estudio es la demanda horaria de energía eléctrica del sistema
peninsular español, entendida como la media, en cada hora, de la demanda real
que publica Red Eléctrica de España a través de su sistema de información
ESIOS. Quedan fuera del análisis los sistemas de los territorios no
peninsulares, que la normativa trata como sistemas eléctricos diferenciados
\parencite{ley2013sectorelectrico}.

La tarea que se aborda es la previsión del día siguiente. A las 12:00 de un
día $D$, hora de cierre del mercado diario, se emiten las previsiones de las
24 horas del día $D+1$. Con esta configuración, la anticipación de las
previsiones oscila entre 12 y 35 horas, o 36 el día del cambio de hora de
octubre, y cada una de ellas solo puede
apoyarse en la información disponible en el momento de su emisión. Esta
restricción condiciona tanto la especificación de los modelos como el diseño
de su evaluación.

La serie presenta tres rasgos que dificultan su modelización. En primer
lugar, superpone varios ciclos estacionales: el diario, de 24 horas, ligado
al ritmo de la actividad humana; el semanal, de 168 horas, que distingue los
días laborables de los fines de semana; y el anual, de unas 8766 horas, la
duración de un año medio de 365,25 días,
asociado a la climatología y a la duración del día. En segundo lugar, está
sometida a efectos de calendario, como los festivos nacionales y autonómicos
o los puentes, que alteran el perfil diario de forma irregular. En tercer
lugar, depende de la temperatura de manera no lineal: la demanda crece tanto
con el frío, por el uso de calefacción, como con el calor, por el de
refrigeración \parencite{hong2016probabilistic}.

Estos rasgos hacen inadecuada la aplicación directa de los modelos clásicos
más extendidos. Un modelo SARIMA incorpora un único periodo estacional, y su
estimación resulta poco práctica cuando ese periodo es largo, como ocurre con
los ciclos semanal y anual de una serie horaria
\parencite{hyndman2021fpp3}. La estacionalidad múltiple exige, por tanto,
recurrir a extensiones específicas: la regresión armónica dinámica, que
representa los ciclos largos mediante términos de Fourier y modela la
dinámica de corto plazo con un proceso ARIMA sobre los errores
\parencite{hyndman2021fpp3}, o los modelos de espacio de estados con
estacionalidad compleja \parencite{delivera2011forecasting}. Los modelos de
aprendizaje automático y profundo, por su parte, pueden aprender estos
patrones directamente de los datos, a cambio de renunciar a una
especificación estadística explícita.

El estudio se delimita, además, mediante tres decisiones cuya justificación
detallada se presenta en el marco metodológico. La primera es el periodo de
estudio, de 2017 a 2024, del que se excluyen los meses de marzo a junio de
2020, en los que las restricciones asociadas a la pandemia de COVID-19
alteraron de forma excepcional el comportamiento de la demanda. La segunda
son las variables explicativas, limitadas a la temperatura horaria ponderada
por la población de cada provincia y al calendario laboral estatal y
autonómico. La tercera, y la de mayores consecuencias, es el tratamiento de
la temperatura. En operación real, a las 12:00 del día $D$ no se conoce la
temperatura del día siguiente, sino una previsión meteorológica sujeta a
error. Este trabajo emplea temperaturas de reanálisis, que reconstruyen las
condiciones observadas, tanto para ajustar como para evaluar los modelos, lo que equivale a suponer
una previsión meteorológica perfecta. Los resultados deben interpretarse, por
tanto, como una cota superior del rendimiento alcanzable en operación, y esta
limitación se retoma al discutir los resultados.

\subsection{Pregunta de investigación y aportación del trabajo}\label{sec:aportacion}

A partir de las limitaciones expuestas en la sección~\ref{sec:carencias}, el
trabajo se plantea la siguiente pregunta de investigación:

\begin{quote}
\itshape
Bajo un protocolo de comparación equitativo, ¿mejoran los modelos de
aprendizaje automático y profundo, de forma estadísticamente significativa, a
un modelo SARIMAX con estacionalidad múltiple en la previsión del día
siguiente de la demanda eléctrica peninsular? ¿En qué horas y tipos de día lo
hacen? ¿Y qué familia de modelos proporciona intervalos de predicción mejor
calibrados?
\end{quote}

Para responderla se comparan cuatro modelos dispuestos en una escala de
complejidad creciente: un modelo ingenuo estacional como referencia mínima, un
modelo SARIMAX con términos de Fourier y variables exógenas, un modelo de
\textit{gradient boosting} \parencite[LightGBM;][]{ke2017lightgbm} y una red
neuronal de previsión multihorizonte \parencite[N-HiTS;][]{challu2023nhits}.
En este trabajo, el término SARIMAX designa una regresión con términos de
Fourier, variables exógenas y errores ARMA, estimada con la clase SARIMAX de
la librería statsmodels; su especificación concreta, sin términos ARIMA
estacionales ni diferenciación, se justifica en la
sección~\ref{sec:esp-sarimax}. Cada familia está representada por un único
modelo, de modo que las conclusiones se refieren a estos cuatro modelos, bajo
este protocolo y en el año de prueba; su extensión al conjunto de cada
familia es una hipótesis que el trabajo no contrasta.

La aportación del trabajo se concreta en tres
elementos. El primero es un protocolo de comparación equitativa, en el que
todos los modelos, incluido el SARIMAX, se evalúan mediante validación con
origen móvil, con la misma información, es decir, la misma demanda pasada y
las mismas variables exógenas, codificada según cada familia, el mismo
horizonte y el mismo presupuesto máximo de configuraciones evaluadas en
validación, fijado y documentado de antemano. El segundo
es una evaluación conjunta de la precisión puntual, con contrastes de
significación de Diebold-Mariano \parencite{diebold1995comparing}, y de la calibración de los intervalos de
predicción, antes y después de recalibrarlos con ACI de forma uniforme para
todos los modelos. El tercero es el desglose de los resultados por hora del
día y tipo de día, que permite identificar en qué situaciones concretas
resulta preferible cada enfoque, en lugar de limitarse a promedios agregados.
El código, las semillas y las particiones temporales se publican para
garantizar la reproducibilidad de los experimentos.

El propósito, por tanto, no es proclamar una familia de modelos como
ganadora. Si los modelos más complejos no mejoran de forma significativa al
SARIMAX, ese resultado constituye un hallazgo del trabajo con implicaciones
prácticas directas, puesto que un modelo más sencillo es también más barato de
mantener y más fácil de interpretar.

\subsection{Estructura del documento}\label{sec:estructura}

El resto del documento se organiza en cinco capítulos. El capítulo 2 formula
el objetivo general y los objetivos específicos del trabajo, junto con los
criterios que permiten considerar cumplido cada uno de ellos.

El capítulo 3 presenta el marco teórico y el estado del arte. Recoge los
fundamentos de las series temporales y su tratamiento cuando presentan
estacionalidad múltiple, los métodos de aprendizaje automático y profundo
aplicados a la previsión, los criterios de evaluación de las previsiones
puntuales y probabilísticas, y la revisión de la literatura sobre previsión de
la demanda eléctrica de la que surgen las carencias señaladas en la
sección~\ref{sec:carencias}.

El capítulo 4 describe el marco metodológico: las fuentes de datos y la
construcción del conjunto de datos, el protocolo de comparación equitativa, la
especificación de cada modelo, el diseño de la capa probabilística y las
métricas y contrastes empleados.

El capítulo 5 expone los resultados: el análisis exploratorio de la serie, la
validación en 2023, en la que se elige la configuración de cada modelo, la
comparación de la precisión puntual en el año de prueba, la evaluación de la calibración de los
intervalos de predicción y el desglose por hora del día y tipo de día,
acompañados de su discusión.

El capítulo 6 recoge las conclusiones, da respuesta a cada uno de los
objetivos específicos y expone las limitaciones del trabajo, entre ellas la
derivada del escenario de previsión meteorológica perfecta, así como las
líneas de trabajo futuro. Por último, los anexos reúnen el material
complementario, como las tablas completas de resultados y el detalle de la
configuración de los modelos.
